\section{Arithmetic resolution towers and nilpotent interactions}\label{sec:arithmetic-towers}

The integral model of \cref{sec:integral-model} has nontrivial exchange maps but identity coarsening functors. We now replace those coarsenings by derived reduction between different coefficient rings. Three questions must be kept separate: whether an exchange is canonically exact, what a specified noninvertible exchange cell leaves behind, and whether forcing that residual to vanish destroys the coefficient category.

\subsection{Derived reduction is a genuine resolution operation}

Fix a finite nonempty set $\mathcal P$ of primes and put
\[
 D=\prod_{p\in\mathcal P}p,\qquad B_n=\mathbb Z/D^n\mathbb Z\quad(n\geq1).
\]
Write $B_{p,n}=\mathbb Z/p^n\mathbb Z$. The Chinese remainder isomorphisms
\[
 B_n\cong\prod_{p\in\mathcal P}B_{p,n}
\]
are compatible with reduction. Use the enhanced category $\mathcal S_n=\operatorname{Perf}(B_n)$ and the coarsening
\[
 C_{m,n}(M)=M\otimes^{\mathbf L}_{B_m}B_n,\qquad m\geq n.
\]
Derived tensor associativity supplies coherent equivalences $C_{l,n}\simeq C_{m,n}C_{l,m}$. These are the data of an enhanced indexed system. Literal equalities are not asserted between arbitrarily chosen derived-functor models. On specified bounded projective complexes the operation is termwise reduction. Preservation of perfect complexes under derived base change is standard \cite{StacksPerfect}; the pointwise cone model also shows exactness directly.

\begin{theorem}[Nonidentity conservative resolution tower]\label{thm:adic-tower}
The functors $C_{m,n}$ are exact and conservative on perfect complexes. They need not be faithful. In particular, for $m>n$, the nonzero endomorphism $D^n:B_m\to B_m$ becomes zero under $C_{m,n}$.
\end{theorem}
\begin{proof}
Exactness follows by reducing bounded projective cone diagrams. For conservativity it suffices, by the Chinese remainder decomposition and composition with reduction to level one, to prove that a perfect complex over $R=\mathbb Z/p^m$ is zero if its derived reduction to $\mathbb F_p$ is zero.

A finite projective module over $R$ is free. Indeed, lift a basis modulo $(p)$ to obtain a surjection $R^r\to P$. It is surjective because its cokernel $Q$ satisfies $Q=pQ$, whence $Q=p^mQ=0$. The surjection splits, and its kernel has zero reduction modulo $p$, so the same argument kills the kernel. A perfect complex therefore has a bounded finite free representative.

If a differential has a unit entry, elementary changes of bases split off the contractible two-term complex $R\xrightarrow{1}R$. The adjacent differentials have zero components into and out of that summand because their composites are zero. Repeating strictly decreases the total rank and terminates in a complex all of whose differential entries lie in $(p)$. Its reduction has zero differentials. If that reduction is acyclic, every remaining rank is zero. The original complex is contractible. Applied to a cone, this also proves that reduction reflects equivalences.

Finally $D^n$ is not zero modulo $D^m$, but is zero modulo $D^n$. Since $\operatorname{Hom}_{D(B_m)}(B_m,B_m)=B_m$, this is genuine loss of morphism information, not just a change of representatives.
\end{proof}

Conservativity detects zero objects and equivalences. It does not recover every morphism, a preferred lift, or the full information in a fine presentation.

\begin{example}[Why ordinary tensor and cohomology truncation are not substitutes]\label{ex:derived-not-naive}
Over $R=\mathbb Z/p^2$, the exact sequence $0\to(p)\to R\to\mathbb F_p\to0$ ceases to be left exact after ordinary tensor with $\mathbb F_p$: the first map becomes zero. The module $\mathbb F_p$ has the infinite periodic free resolution with every differential multiplication by $p$. Tensoring that resolution with $\mathbb F_p$ gives zero differentials, hence $\operatorname{Tor}_i^R(\mathbb F_p,\mathbb F_p)=\mathbb F_p$ for every $i\geq0$.

In contrast, $C_R(p)=[R\xrightarrow p R]$ is perfect, but has both $H^{-1}=\mathbb F_p$ and $H^0=\mathbb F_p$. It is not the module $\mathbb F_p$ placed in degree zero. We retain the actual two-term complex throughout the tower. Derived tensor is computed on these projective representatives, not on a truncated cohomology surrogate \cite{StacksDerivedTensor}.
\end{example}

\subsection{Coherent cells and their actual source of residuals}

For a fixed perfect integral complex $L$, put
\[
 A_{L,n}(M)=(L\otimes^{\mathbf L}_{\mathbb Z}B_n)\otimes^{\mathbf L}_{B_n}M.
\]
Base change gives a natural equivalence
\[
 \theta_{L;m,n}:A_{L,n}C_{m,n}\simeq C_{m,n}A_{L,m}.
\]
Thus the canonical tensor/base-change exchange already has zero residual. The tower alone does not force a defect. To test a noninvertible exchange, we explicitly add a positive integer label $b$ and use the cell
\[
 \eta_{L,b;m,n}=b^{m-n}\theta_{L;m,n}.
\]
For scalar unit templates take $L=\mathbb Z$, so that the transformation functors are identities but the coarsening functors are not.

\begin{theorem}[Coherent weighted exchange and residual transport]\label{thm:weighted-adic-exchange}
The cells $\eta_{L,b;m,n}$ are natural and coherent with composition of resolution changes. Their residual at $M\in\mathcal S_m$ is
\[
 \Cofib(\eta_{L,b;m,n,M})\simeq
 A_{L,n}C_{m,n}(M)\otimes^{\mathbf L}_{B_n}C_{B_n}(b^{m-n}).
\]
For a fixed scalar $a$, subsequent base change transports its cone exactly:
\[
 C_{n,k}\bigl(N\otimes^{\mathbf L}_{B_n}C_{B_n}(a)\bigr)
 \simeq C_{n,k}(N)\otimes^{\mathbf L}_{B_k}C_{B_k}(a).
\]
At a common target level $n$, two independent scalar events of positive lengths $d,e$ and labels $a,b$ have interaction
\[
 N\otimes^{\mathbf L}_{B_n}K_{B_n}(a^d,b^e),\qquad
 K_R(u,v)=C_R(u)\otimes^{\mathbf L}_R C_R(v).
\]
These formulas also hold with compatible tensor templates absorbed into $N$.
\end{theorem}
\begin{proof}
Naturality follows from central scalar multiplication. For $l\geq m\geq n$, the two-step cell has scalar $b^{m-n}b^{l-m}=b^{l-n}$. The associativity coherence of derived tensor identifies its underlying $\theta$ with $\theta_{L;l,n}$. Identity intervals have scalar one. Tensoring transformation templates gives their usual associative action, with scalar weights multiplied; no additional cells are inferred without these specified data.

After using $\theta$ to identify the target, the exchange is multiplication by $b^{m-n}$ on $A_{L,n}C_{m,n}M$. Exact tensor identifies its cone with the first formula. Associativity of derived tensor proves the second. The two-variable interaction formula follows by taking the two cofibers, as in \cref{thm:tensor-cube}. For a pure adjacent-swap word, tensor templates keep their relative order and all coarsenings combine into the same final reduction. Identifying those endpoints by the supplied tensor associativities leaves precisely the products of the scalar cells. This constructs the corresponding coherent scalar event diagram at the target level.
\end{proof}

\begin{remark}[Specified cells, not a claimed arithmetic correspondence]
Nonzero residuals in this model come from the chosen noninvertible scalar cells. They are not defects of the canonical base-change equivalence, and they are not identified with a ramification or automorphy obstruction. The construction supplies a reproducible nonidentity resolution test bed with explicit semantics.
\end{remark}

\subsection{The finite-level Koszul calculation}

\begin{theorem}[Prime-power interaction cohomology]\label{thm:prime-power-koszul}
Let $R=\mathbb Z/p^n$, let $a_1,\ldots,a_t$ be positive integers with $t\geq1$, and put
\[
 h=\min\{n,v_p(a_1),\ldots,v_p(a_t)\},\qquad Q_h=R/(p^h).
\]
Here $Q_0=0$ and $Q_n=R$. There is a generally noncanonical chain isomorphism
\[
 K_R(a_1,\ldots,a_t)\cong
 \bigoplus_{j=0}^{t-1}C_R(p^h)[j]^{\oplus\binom{t-1}{j}}.
\]
Consequently
\[
 \boxed{H^{-j}K_R(a_1,\ldots,a_t)\cong Q_h^{\oplus\binom tj}
 \quad(0\leq j\leq t),}
\]
and all other cohomology groups vanish. The complex is zero in $D(R)$ exactly when at least one label is a unit modulo $p$.
\end{theorem}
\begin{proof}
If $h=n$, all entries are zero in $R$. Otherwise choose an entry of minimal valuation, scale its unit part away, and subtract suitable multiples of it from the other entries. These operations form an invertible change of the coefficient row and produce $(p^h,0,\ldots,0)$. The exterior powers of the basis change give a Koszul chain isomorphism, with the convention of \cref{thm:integral-gcd} and \cite{StacksKoszul}.

Each zero entry contributes $R\oplus R[1]$. This proves the complex decomposition. Multiplication by $p^h$ on $R$ has cokernel $Q_h$ and kernel $p^{n-h}R$, the latter isomorphic to $Q_h$ by multiplication by $p^{n-h}$. Hence $C_R(p^h)$ has $Q_h$ in degrees $-1$ and $0$. Counting the shifted copies and using Pascal's identity gives the formula. The unit case is contractible, and otherwise $H^0=Q_h\ne0$.
\end{proof}

For two events the cohomology ranks over $Q_h$ are $1,2,1$ in degrees $-2,-1,0$, rather than the $0,1,1$ pattern over the integral or complete domain. These extra classes are real finite-level cohomology. They may not be discarded while retaining the same complex. In particular, a cohomology group appearing here is not automatically a perfect $R$-module.

\begin{theorem}[Nilpotent interaction forces excisive collapse]\label{thm:nilpotent-collapse}
For any commutative ring $R$ and nilpotent $x\in R$, the thick subcategory generated by $C_R(x)$ is all of $\operatorname{Perf}(R)$. Therefore, over $R=\mathbb Z/p^n$, forcing the interaction of two nonunit scalar labels to vanish in an exact quotient forces that quotient to be zero.
\end{theorem}
\begin{proof}
Choose $N$ with $x^N=0$. For composable scalar multiplications, the octahedral cofiber relation builds $C_R(x^k)$ from $C_R(x^{k-1})$ and $C_R(x)$. Induction puts $C_R(x^N)$ in $\operatorname{Thick}(C_R(x))$. But
\[
 C_R(x^N)=C_R(0)\cong R\oplus R[1].
\]
Retract closure therefore includes $R$, and $R$ thickly generates $\operatorname{Perf}(R)$. This cone-of-a-product mechanism is also part of the standard Koszul calculus \cite{StacksKoszul}.

For two nonunit labels, \cref{thm:prime-power-koszul} makes their interaction a sum of $C_R(p^h)$ and its shift, with $1\leq h\leq n$. Thus it contains $C_R(p^h)$ as a retract. This element is nilpotent. Any exact functor annihilating the interaction annihilates the thick subcategory it generates, hence every perfect object.
\end{proof}

This is stronger than the zero-label counterexample over $\mathbb Z$. At finite $p$-power resolution, even nonzero labels such as $p$ and $p^2$ can force total collapse. Localizing $\mathbb Z/p^n$ at $p$ gives the zero ring. Noncollapse over the integral coefficients does not survive this base change automatically.

\subsection{Several primes give a nonzero, but selective, quotient}

Specify finitely many positive scalar labels and a collection $\mathcal I$ of allowed independent pairs. Include their unit-seed interaction diagrams as the designated tests. This hypothesis ensures that the minimal kernel below is not being inferred from templates that accidentally annihilate an interaction already. Put
\[
 S=\{p\in\mathcal P:\text{some }(a,b)\in\mathcal I
                    \text{ satisfies }p\mid a\text{ and }p\mid b\},
 \qquad D_0=\prod_{p\in\mathcal P\setminus S}p.
\]
Let $\mathcal J_n$ be the tensor-invariant thick closure of those interactions in $\operatorname{Perf}(B_n)$. Positive lengths in \cref{thm:weighted-adic-exchange} do not change $S$.

\begin{theorem}[Prime-component quotient and tower compatibility]\label{thm:finite-prime-quotient}
Under the Chinese remainder equivalence,
\[
 \boxed{\mathcal J_n=
 \prod_{p\in S}\operatorname{Perf}(\mathbb Z/p^n)
 \times\prod_{p\in\mathcal P\setminus S}0.}
\]
In particular there is an equivalence of stable categories
\[
 \operatorname{Perf}(B_n)/\mathcal J_n
 \simeq \prod_{p\in\mathcal P\setminus S}\operatorname{Perf}(\mathbb Z/p^n)
 \simeq\operatorname{Perf}(\mathbb Z/D_0^n).
\]
The equivalence is realized by projection to the retained prime components. Its kernel is constructor-invariant and preserved by every coarsening. If $q_n$ denotes projection, then
\[
 q_nC_{m,n}\simeq C^0_{m,n}q_m
\]
with the supplied base-change coherence. The quotient is nonzero exactly when $S\ne\mathcal P$. When $S=\mathcal P$, use the zero category and the zero ring $\mathbb Z/1$.
\end{theorem}
\begin{proof}
At a prime outside $S$, each designated interaction has a unit label and is acyclic. At a prime in $S$, at least one pair has both labels nonunit. Tensor by the corresponding Chinese remainder idempotent to isolate that component. The nilpotent-collapse theorem then puts its unit in $\mathcal J_n$. Its whole perfect category follows. These two observations prove equality of the kernel with the displayed product subcategory.

Projection to the complementary factors is exact and has an exact fully faithful section. Every object is a direct sum of its selected and retained components. In the quotient the selected summand is zero, so that object becomes isomorphic to its retained summand. Projection and section therefore induce inverse equivalences. This proves the whole quotient assertion directly, without an unproved identification of an arbitrary Verdier quotient with localized perfect complexes.

Tensor products and reduction respect the Chinese remainder idempotents. Projection consequently commutes with reduction, and the same tensor associativities give coherence for composite reductions. All prescribed square interactions vanish after projection, so the earlier excisive reconstruction results apply at every level.
\end{proof}

This repair does not retain the selected prime's finer torsion while deleting only one common exponent. It removes the entire selected prime component. Thus its minimality is a statement about exact excision, not a guarantee that it preserves all arithmetic information desired by a later application.
